
\documentclass{article}
\usepackage{colortbl}
\usepackage{makecell}
\usepackage{multirow}
\usepackage{supertabular}

\begin{document}

\newcounter{utterance}

\centering \large Interaction Transcript for game `adventuregame', experiment `home\_deliver\_three\_basic\_hard\_invlimittwo', episode 7 with reka{-}flash{-}3.1{-}t0.0.
\vspace{24pt}

{ \footnotesize  \setcounter{utterance}{1}
\setlength{\tabcolsep}{0pt}
\begin{supertabular}{c@{$\;$}|p{.15\linewidth}@{}p{.15\linewidth}p{.15\linewidth}p{.15\linewidth}p{.15\linewidth}p{.15\linewidth}}
    \# & \multicolumn{2}{c}{Player} && \multicolumn{2}{c}{Game Master} \\
    \hline

    \theutterance \stepcounter{utterance}  
    & & & \multicolumn{4}{p{0.6\linewidth}}{
        \cellcolor[rgb]{0.9,0.9,0.9}{
            \makecell[{{p{\linewidth}}}]{
                \texttt{\tiny{[P1$\langle$GM]}}
                \texttt{You are playing a text adventure game. I will describe what you can perceive in the game. You write the single action you want to take in the game starting with >. Only reply with an action.} \\
\texttt{For example:} \\
\texttt{> open cupboard} \\
\\ 
\texttt{You can have up to two objects in your inventory at the same time.} \\
\\ 
\texttt{Your goal for this game is: Put the potted plant in the freezer, the plate on the table and the mop in the refrigerator.} \\
\\ 
\texttt{Once you have achieved your goal, write "> done" to end the game.} \\
\\ 
\texttt{You are in a bedroom now. There are a pillow, a wardrobe and a bed. The pillow is in the wardrobe. The wardrobe is closed. There is a passage to a hallway here.} \\
            }
        }
    }
     \\ \\

    \theutterance \stepcounter{utterance}  
    & \multicolumn{4}{p{0.6\linewidth}}{
        \cellcolor[rgb]{0.561,0.737,0.561}{
            \makecell[{{p{\linewidth}}}]{
                \texttt{\tiny{[P1$\rangle$GM]}}
                \texttt{<reasoning>} \\
\texttt{Let's see. The goal is to put the potted plant in the freezer, the plate on the table, and the mop in the refrigerator. But right now, I'm in the bedroom. The initial description says there's a pillow, wardrobe, bed, and a passage to the hallway. The pillow is in the wardrobe, which is closed. So first, maybe I need to open the wardrobe to get the pillow out? But why would I need the pillow? The goals don't mention the pillow. Maybe the wardrobe has other items? The problem is, the user hasn't mentioned a potted plant, plate, or mop in the bedroom. So maybe those items are in other rooms, like the hallway or kitchen? The passage to the hallway is here, so perhaps I need to go there first.} \\
\\ 
\texttt{Wait, the initial description only lists the bedroom's contents. The other items (potted plant, plate, mop) must be in other rooms. So the first step is to go to the hallway via the passage. Let me check the possible actions. The user can move through passages, open/close objects, take items, put items down, etc. So the first action should be moving to the hallway. Let me confirm: "There is a passage to a hallway here." So I can go to the hallway. Let me do that. So the action would be "> go to hallway". Then, in the hallway, maybe there are more rooms or items. But the user hasn't described the hallway yet. Wait, the initial description is only the bedroom. So when I move to the hallway, the next description will come. But since the user hasn't provided that yet, maybe I need to proceed based on the given info. Wait, no—the user is going to describe the environment as I take actions. So when I go to the hallway, the next step will be the description of the hallway. But since the user hasn't provided that yet, perhaps I need to make an assumption. Wait, no. The user is going to provide the next description after each action. So in the current state, the user is in the bedroom. The first action is to go to the hallway. Let me proceed with that. So the first action is "> go to hallway". Then, the user will describe the hallway's contents. From there, maybe there's a freezer, refrigerator, table, etc. But the problem is, the user hasn't} \\
            }
        }
    }
    & & \\ \\

    \theutterance \stepcounter{utterance}  
    & & & \multicolumn{2}{p{0.3\linewidth}}{
        \cellcolor[rgb]{0.9,0.9,0.9}{
            \makecell[{{p{\linewidth}}}]{
                \texttt{\tiny{[GM$|$GM]}}
                \texttt{Round: 0} \\
            }
        }
    }
    & & \\ \\

    \theutterance \stepcounter{utterance}  
    & & & \multicolumn{2}{p{0.3\linewidth}}{
        \cellcolor[rgb]{0.9,0.9,0.9}{
            \makecell[{{p{\linewidth}}}]{
                \texttt{\tiny{[GM$|$GM]}}
                \texttt{command\_tag\_missing} \\
            }
        }
    }
    & & \\ \\

    \theutterance \stepcounter{utterance}  
    & & & \multicolumn{2}{p{0.3\linewidth}}{
        \cellcolor[rgb]{0.9,0.9,0.9}{
            \makecell[{{p{\linewidth}}}]{
                \texttt{\tiny{[GM$|$GM]}}
                {'goal_states_achieved': [], 'game_successfully_finished': False}
            }
        }
    }
    & & \\ \\

    \theutterance \stepcounter{utterance}  
    & & & \multicolumn{2}{p{0.3\linewidth}}{
        \cellcolor[rgb]{0.9,0.9,0.9}{
            \makecell[{{p{\linewidth}}}]{
                \texttt{\tiny{[GM$|$GM]}}
                \texttt{command\_tag\_missing} \\
            }
        }
    }
    & & \\ \\

\end{supertabular}
}

\end{document}
